\section{Introduction}

Document understanding is a core capability in modern information systems, supporting the extraction and structuring of knowledge from visually rich and layout-intensive documents such as financial reports, scientific articles, contracts, and invoices. Traditional OCR systems~\cite{Tesseract,EasyOCR,PaddleOCR} mainly focus on plain text transcription and rely on multi-stage pipelines with handcrafted rules for layout parsing and downstream information extraction. While effective for simple scenarios, these approaches often struggle with complex layouts, diverse document formats, and real-world production requirements.

Recent multimodal large language models (MLLMs)~\cite{Qwen3-VL,glm-4.5v,seed1_5vl} unify visual perception and language understanding within a single framework and significantly improve document understanding performance. However, their large model size and autoregressive decoding paradigm lead to high computational cost, slow inference, and substantial memory consumption, which makes large-scale deployment under high-concurrency or edge environments challenging.

In practical production systems, document intelligence solutions must simultaneously provide: (1) strong performance on complex content such as tables, formulas, code, and seals, (2) high-throughput and low-latency inference, and (3) flexible integration and domain adaptation. GLM-OCR is developed to address these system-level requirements within a unified multimodal framework.

\textbf{GLM-OCR} is a lightweight multimodal OCR model for comprehensive document understanding. Built on the GLM-V encoder–decoder framework~\cite{glm-4.5v}, it combines a 0.4B-scale CogViT visual encoder trained in large-scale image–text data, a lightweight cross-modal connector and a 0.5B-scale GLM language decoder~\cite{glm2024chatglm}. The entire model contains only 0.9B parameters, enabling high-throughput and low-latency inference while maintaining strong recognition performance.

Beyond architectural optimization, GLM-OCR also considers the mismatch between conventional autoregressive generation and the characteristics of OCR tasks. OCR is inherently a deterministic task with strong local dependencies and explicit structural supervision, where strictly autoregressive token-by-token decoding is inefficient. Therefore, we introduce Multi-Token Prediction (MTP)~\cite{deepseekv3} into both training and inference. MTP enables the simultaneous prediction of multiple tokens, substantially improving training efficiency and decoding throughput while preserving recognition accuracy, and is particularly advantageous for long structured outputs such as tables. To control the additional memory overhead introduced by MTP, we further adopt a parameter-sharing scheme across the draft models, which substantially reduces the additional GPU memory overhead~\cite{glm5}. In practice, GLM-OCR is trained to predict ten tokens per step and generates 5.2 tokens per decoding step on average at inference time, bringing approximately 50\% throughput improvement.

At the system level, GLM-OCR adopts a two-stage pipeline consisting of layout analysis and parallel content recognition. The layout stage is powered by PP-DocLayout-V3~\cite{PaddleOCR-VL-1.5}, which detects structured regions and enables parallel recognition across different document areas. This design improves both robustness and processing efficiency for complex real-world documents.

\textbf{Results.} Figure~\ref{fig:eval_comparison} shows that GLM-OCR achieves 94.6 on OmniDocBench v1.5~\cite{OmniDocBench}, ranking first among all evaluated models despite its compact 0.9B size. Besides, GLM-OCR delivers strong performance across text recognition, formula recognition, table parsing, and key information extraction, reaching 94.0 on OCRBench (Text)~\cite{OCRBench} and 96.5 on UniMERNet~\cite{UniMERNet}, and achieving 85.2 on PubTabNet~\cite{PubTabNet} and 86.0 on TEDS~\footnote{https://modelscope.cn/datasets/jockerK/TEDS\_TEST}. It also performs competitively on information extraction benchmarks such as Nanonets-KIE and Handwritten-Forms~\cite{IDPLeaderboard}, with performance comparable to significantly larger general multimodal models.

In addition to public benchmarks, we evaluate GLM-OCR on six high-frequency real-world scenarios, including code document parsing, natural-scene table recognition, handwritten text recognition, multilingual OCR~\footnote{Chinese, English, French, Spanish, Russian, German, Japanese, and Korean}, seal recognition, and receipt KIE. GLM-OCR consistently delivers strong results across all settings, achieving 91.5 on real-world table recognition, 90.5 on seal recognition, and 94.5 on receipt KIE. These results indicate that GLM-OCR generalizes beyond curated benchmarks and remains effective under practical production conditions.

\textbf{Deployment and Finetuning.} GLM-OCR supports efficient inference with modern serving frameworks such as vLLM~\cite{vllm}~\footnote{https://github.com/vllm-project/vllm}, SGLang~\footnote{https://github.com/sgl-project/sglang}, and Ollama~\footnote{https://github.com/ollama/ollama}, enabling deployment in both large-scale and resource-constrained edge scenarios. It also provides full finetuning support through LLaMA-Factory~\footnote{https://github.com/hiyouga/LlamaFactory}, enabling rapid adaptation to domain-specific document understanding tasks~\footnote{Tutorial: https://github.com/zai-org/GLM-OCR/blob/main/examples/finetune/README.md}. 


% --- 论文风味重一些 ---

% Document understanding plays a central role in modern information processing systems by enabling the extraction, organization, and structuring of knowledge from visually rich and layout-intensive documents, such as financial reports, scientific articles, contracts, and invoices. Traditional optical character recognition (OCR) systems~\cite{Tesseract,EasyOCR,PaddleOCR} mainly focus on plain text transcription and typically rely on multi-stage pipelines with handcrafted rules for layout analysis and downstream information extraction. Although effective for basic text recognition, these pipeline-based approaches often lack robustness and flexibility when dealing with complex layouts and diverse document formats. 

% Recent multimodal large language models (MLLMs)~\cite{Qwen3-VL,glm-4.5v,seed1_5vl} have significantly advanced unified document understanding by integrating visual perception and language reasoning into a single framework. However, despite their strong performance, these models usually come with high computational cost, slow inference speed, and limited deployability under high-concurrency or resource-constrained settings.

% In real-world applications, document intelligence systems~\cite{zhang2024document,ke2025large} must simultaneously satisfy several system-level requirements: (1) robust performance across diverse and complex layouts, including tables, formulas, code snippets, and seals; (2) efficient inference with low latency and low memory footprint; and (3) seamless integration into production pipelines. Existing end-to-end MLLM-based OCR solutions often prioritize benchmark accuracy while overlooking these practical constraints, making them difficult to deploy at scale.

% To address these challenges, we present \textbf{GLM-OCR}, a lightweight multimodal OCR model for comprehensive document understanding. Built upon the GLM-V encoder–decoder framework~\cite{glm-4.5v}, GLM-OCR integrates a visually powerful CogViT encoder pretrained on large-scale image–text data, a lightweight cross-modal connector, and a 0.5B-scale GLM language decoder~\cite{glm2024chatglm}. This design results in a compact model with only 0.9B parameters, enabling high-throughput and low-latency inference while maintaining strong recognition performance.

% Beyond architectural optimization, GLM-OCR revisits the generation paradigm of language models for OCR. Given the deterministic nature of OCR, together with its strong local dependencies and explicit structural supervision, strictly autoregressive token-by-token decoding is often unnecessary and inefficient. We therefore introduce Multi-Token Prediction (MTP)~\cite{deepseekv3} into both training and inference. MTP enables the simultaneous prediction of multiple tokens, substantially improving training efficiency and decoding throughput while preserving recognition accuracy, and is particularly advantageous for long structured outputs such as tables. To control the additional memory overhead introduced by MTP, we further adopt a parameter-sharing scheme across the draft models~\cite{glm5}, significantly reducing the GPU memory footprint. In practice, the model is trained to predict 10 tokens per step and generates 5.2 tokens per decoding step on average during inference, yielding an approximately 50\% throughput improvement.

% From a system perspective, GLM-OCR follows a two-stage pipeline consisting of layout analysis and parallel content recognition. The layout stage is powered by PP-DocLayout-V3~\cite{PaddleOCR-VL-1.5}, which enables structured region detection and allows the recognition model to process different regions in parallel, thereby improving both robustness and efficiency on complex real-world documents.

% \textbf{Results.} Figure~\ref{fig:eval_comparison} shows that GLM-OCR achieves 94.6 on OmniDocBench v1.5~\cite{OmniDocBench}, ranking first overall. Despite having only 0.9B parameters, it consistently outperforms existing OCR-specialized systems across multiple benchmarks, including text recognition, formula recognition, table structure parsing, and key information extraction, while achieving performance comparable to state-of-the-art general multimodal models. Specifically, GLM-OCR obtains 94.0 on OCRBench (Text)~\cite{OCRBench} and 96.5 on UniMERNet~\cite{UniMERNet}, and delivers strong table parsing performance with 85.2 on PubTabNet~\cite{PubTabNet} and 86.0 on TEDS~\footnote{https://modelscope.cn/datasets/jockerK/TEDS\_TEST}. It also achieves competitive results on information extraction benchmarks such as Nanonets-KIE and Handwritten-Forms~\cite{IDPLeaderboard}.

% Beyond public benchmarks, we further conduct internal evaluations on six representative real-world scenarios, including code document parsing, complex table recognition in natural images, handwritten text recognition, multilingual OCR, seal recognition, and receipt KIE. GLM-OCR demonstrates clear and consistent advantages across all settings. In particular, it achieves 91.5 on real-world table recognition, 90.5 on seal recognition, and 94.5 on receipt KIE, substantially surpassing other OCR-specific models and in several cases approaching the performance of much larger general vision-language models. These results verify that GLM-OCR not only excels on standardized benchmarks but also generalizes effectively to diverse practical applications.

% \textbf{Deployment and Finetuning.} GLM-OCR supports efficient inference with modern serving frameworks such as vLLM~\cite{vllm}~\footnote{https://github.com/vllm-project/vllm}, SGLang~\footnote{https://github.com/sgl-project/sglang}, and Ollama~\footnote{https://github.com/ollama/ollama}, enabling deployment in both large-scale and resource-constrained edge scenarios. It also provides full finetuning support through LLaMA-Factory~\footnote{https://github.com/hiyouga/LlamaFactory}, allowing users to adapt the model to domain-specific document understanding tasks~\footnote{Tutorial: https://github.com/zai-org/GLM-OCR/blob/main/examples/finetune/README.md}.